\documentclass[../DoAn.tex]{subfiles}
\begin{document}

% Nguồn số liệu của chương: repo SUMO (CLAUDE.md, pyproject.toml, docs/PROGRESS.md,
% docs/research/phase3_box_capacity.md, experiments/phase3/results_r21.csv,
% experiments/operator/task4b/b0_table.csv, src/thesis_sumo/**; cập nhật 2026-10-27: docs/research/phase3*.md,
% report/2026-10-27/data/campaign_specs.csv, experiments/campaigns/{P1,P2}/loop/iter_*.json,
% experiments/llm_logs/{P1,P2}/validity_log.csv, experiments/campaigns/REF_run.log). Các mục [TODO] cần bổ sung trước khi nộp.

Chương này trình bày phần hiện thực của phương pháp ở Chương~\ref{chapter:Methodology}: kiến trúc phần mềm, các gói và lớp chính, công cụ sử dụng, kiểm thử, các thực nghiệm đã tiến hành và cách tái lập. Tại thời điểm viết (27/10/2026), toàn bộ mã của quy trình đã được hiện thực, từ mạng mô phỏng, bộ điều khiển TraCI, các chỉ số, khung đánh giá đến vòng lặp LLM và phương án Optuna; chiến dịch Optuna B2 đã hoàn thành, hai chiến dịch LLM P1 và P2 mới có kết quả của 24/36 lần chạy (2/3 vòng), còn phương án H và đánh giá cuối chưa chạy (Bảng~\ref{table:3.1}). Mọi kết quả định lượng trong chương được tính trên mạng và dữ liệu đại diện (\texttt{LITERATURE\_PROXY}) xây dựng từ tài liệu, không phải số liệu đo tại nút giao; mọi đặc tả P1, P2 được sinh bởi mô hình thay thế Claude Sonnet chứ không phải LLM cục bộ (Mục~\ref{section:3.8}).

\section{Thiết kế kiến trúc}

\subsection{Lựa chọn kiến trúc phần mềm}

Hệ thống được tổ chức theo kiểu đường ống (pipeline) gồm các bước tách rời, mỗi bước đọc tệp đầu vào và ghi tệp đầu ra ở dạng văn bản (YAML, JSON, XML, CSV). Lựa chọn này xuất phát từ ba yêu cầu. Thứ nhất là tính tái lập: mạng lưới được sinh lại từ dữ liệu gốc bằng script và các tệp vá, mọi lần chạy dùng seed cố định và ghi một tệp \texttt{manifest.json}. Thứ hai là khả năng thay thế nguồn đặc tả luật: B0, B1, Optuna và LLM đều chỉ tạo ra một đặc tả JSON cùng lược đồ, phần còn lại của đường ống không phân biệt nguồn gốc của đặc tả. Thứ ba là khả năng dùng chung một bộ hàm chỉ số cho dữ liệu mô phỏng và dữ liệu quan sát, nhờ chuyển đầu ra SUMO về cùng lược đồ với dữ liệu quan sát.

Mã nguồn là một gói Python duy nhất \texttt{thesis\_sumo}, chia thành các gói con theo trách nhiệm (Bảng~\ref{table:4.1}). Ngoài các gói con, mô-đun \texttt{evaluate} chứa khung đánh giá dùng chung cho mọi phương án.

\begin{table}[H]
\centering
\small
\begin{tabular}{|p{2.4cm}|p{8.3cm}|p{2.3cm}|}
\hline
\textbf{Gói} & \textbf{Trách nhiệm} & \textbf{Trạng thái} \\ \hline
\texttt{config}, \texttt{schema\_\allowbreak observed} & Đọc tệp cấu hình hiện trường, thí nghiệm và LLM; báo lỗi khi trường bắt buộc còn trống; định nghĩa các bảng và cột của lược đồ dữ liệu quan sát & Đã triển khai \\ \hline
\texttt{net} & Sinh mạng từ OpenStreetMap bằng \texttt{netconvert}, áp tệp vá plain XML (bề rộng làn, kết nối rẽ, chương trình đèn), xuất bảng cạnh và sơ đồ mạng & Đã triển khai \\ \hline
\texttt{demand} & Sinh tuyến ứng viên và tệp tuyến theo số đếm rẽ 15 phút bằng công cụ \texttt{routeSampler} của SUMO & Đã triển khai \\ \hline
\texttt{proxy} & Sinh năm ngày dữ liệu đại diện PROXY\_D1--PROXY\_D5 theo từng sự kiện phương tiện từ bảng giá trị mục tiêu & Đã triển khai \\ \hline
\texttt{rules} & Lược đồ đặc tả luật, bộ kiểm tra hợp lệ, bộ biên dịch & Đã triển khai \\ \hline
\texttt{runtime} & Bộ điều khiển TraCI và năm luật có điều kiện & Đã triển khai \\ \hline
\texttt{sim} & Ghi cấu hình SUMO, chạy mô phỏng, ghi manifest, chuyển đầu ra sang lược đồ quan sát & Đã triển khai \\ \hline
\texttt{metrics} & Các chỉ số M1--M10, cổng G0, hàm khoảng cách và điểm tổng & Đã triển khai \\ \hline
\texttt{evaluate} & Bộ đếm ngân sách theo chiến dịch, đánh giá một đặc tả trên các ngày hiệu chỉnh, chọn đặc tả tốt nhất & Đã triển khai \\ \hline
\texttt{llm} & Gọi LLM (backend tương thích OpenAI và backend tệp cho mô hình thay thế), prompt \texttt{v001}, tạo phản hồi, vòng lặp tinh chỉnh P1/P2 & Đã triển khai \\ \hline
\texttt{baselines} & Phương án đối sánh B2 và H bằng Optuna & Đã triển khai \\ \hline
\end{tabular}
\caption{Các gói con của \texttt{thesis\_sumo}}
\label{table:4.1}
\end{table}

\begin{sloppypar}
Điểm cốt lõi của kiến trúc là tách hai giai đoạn theo thời gian. Giai đoạn ngoại tuyến (offline) gồm việc sinh hoặc soạn đặc tả luật, kiểm tra hợp lệ và biên dịch; LLM (hoặc mô hình thay thế trong các chiến dịch hiện tại) chỉ tham gia ở giai đoạn này và chỉ tạo ra tệp JSON. Giai đoạn trực tuyến (online) là một lần chạy SUMO: tham số tĩnh đã được SUMO nạp từ tệp \texttt{vtypes.add.xml} khi khởi tạo, còn bộ điều khiển TraCI đọc \texttt{traci\_rules.json} và đánh giá các luật có điều kiện tại mỗi bước mô phỏng. Không có lời gọi LLM nào trong vòng lặp mô phỏng; LLM không điều khiển phương tiện theo thời gian thực.
\end{sloppypar}

Mỗi lần chạy được điều phối bởi một trình chạy (runner) duy nhất. Trình chạy xác định đầu vào theo quy ước tên (mạng theo tệp cấu hình hiện trường, tệp tuyến theo ngày, đặc tả đã biên dịch theo mã đặc tả), khởi động SUMO qua thư viện \texttt{libsumo} (chạy hàng loạt) hoặc \texttt{traci} (khi cần giao diện đồ họa \texttt{sumo-gui}), gắn bộ điều khiển TraCI nếu đặc tả có luật, rồi ghi tệp manifest và chuyển đầu ra sang lược đồ quan sát. Nhiều lần chạy (nhiều seed) được thực hiện song song, mỗi tiến trình một phiên SUMO.

\subsection{Thiết kế tổng quan}

Luồng dữ liệu của hệ thống được tóm tắt trong Bảng~\ref{table:4.2}. Các bước 1--3 chạy một lần cho mỗi cấu hình hiện trường hoặc mỗi ngày dữ liệu; các bước 4--8 lặp lại cho mỗi đặc tả luật; bước 9 đóng vòng lặp bằng cách đưa kết quả về cho bộ sinh đặc tả.

\begin{table}[H]
\centering
\small
\begin{tabular}{|p{0.6cm}|p{3.4cm}|p{3.6cm}|p{4.9cm}|}
\hline
\textbf{\#} & \textbf{Bước (gói)} & \textbf{Đầu vào} & \textbf{Đầu ra} \\ \hline
1 & Dựng mạng (\texttt{net}) & Bản xuất OSM, tệp cấu hình hiện trường & Mạng cơ sở, mạng hiện trường, tệp vá, bảng cạnh, sơ đồ mạng \\ \hline
2 & Dữ liệu mục tiêu (\texttt{proxy}) & Bảng giá trị mục tiêu, cấu hình hiện trường & Các bảng CSV theo lược đồ quan sát của từng ngày \\ \hline
3 & Demand (\texttt{demand}) & Mạng, số đếm rẽ 15 phút & Tệp tuyến của từng ngày \\ \hline
4 & Kiểm tra hợp lệ (\texttt{rules}) & Đặc tả JSON, miền tham số, cấu hình hiện trường, mạng & Danh sách lỗi có mã \\ \hline
5 & Biên dịch (\texttt{rules}) & Đặc tả hợp lệ & \texttt{vtypes.add.xml}, \texttt{traci\_\allowbreak rules.json} (nếu có luật) \\ \hline
6 & Chạy mô phỏng (\texttt{sim}, \texttt{runtime}) & Mạng, tệp tuyến, đặc tả đã biên dịch, cấu hình thí nghiệm, seed & Đầu ra SUMO, \texttt{sumo.log}, \texttt{interventions.csv}, \texttt{manifest.json} \\ \hline
7 & Chuyển lược đồ (\texttt{sim}) & Đầu ra SUMO & Các bảng CSV theo lược đồ quan sát \\ \hline
8 & Chỉ số và điểm (\texttt{metrics}, \texttt{evaluate}) & Bảng CSV mô phỏng và quan sát của các ngày hiệu chỉnh & M1--M10, G0, $R$, $S$, tính đủ điều kiện \\ \hline
9 & Sinh đặc tả kế tiếp (\texttt{llm}, \texttt{baselines}) & Kết quả đánh giá, ngân sách còn lại & Phản hồi cho LLM hoặc giá trị mục tiêu cho Optuna; đặc tả mới \\ \hline
\end{tabular}
\caption{Luồng dữ liệu giữa các gói}
\label{table:4.2}
\end{table}

\begin{sloppypar}
Phụ thuộc giữa các gói đi theo một chiều. Gói \texttt{config} và \texttt{schema\_observed} là nền tảng, được mọi gói khác sử dụng. Gói \texttt{rules} dùng chung một mô-đun xác định kết nối đi sai làn với gói \texttt{runtime}, để bộ kiểm tra hợp lệ và bộ điều khiển có cùng cách hiểu về các kết nối có thể dùng. Gói \texttt{sim} gọi \texttt{runtime} khi đặc tả có luật TraCI. Gói \texttt{proxy} sinh dữ liệu mục tiêu độc lập với mọi đầu ra SUMO, nhưng dùng lại các định nghĩa hình học và lược đồ của \texttt{sim} để hai phía có cùng định nghĩa bảng. Gói \texttt{metrics} chỉ đọc các bảng CSV theo lược đồ quan sát, nên không phụ thuộc vào SUMO.
\end{sloppypar}

\subsection{Thiết kế chi tiết gói}
\label{subsection:package_design}

Phần này mô tả các thành phần chính như đã hiện thực; chức năng của bộ kiểm tra hợp lệ, bộ biên dịch và bộ điều khiển đã được trình bày ở Mục~\ref{section:3.4} và Mục~\ref{section:3.5}.

\textbf{Gói \texttt{config}.} Các hàm đọc cấu hình hiện trường, cấu hình thí nghiệm và cấu hình LLM trả về đối tượng cấu hình; khi một trường bắt buộc có giá trị rỗng, chương trình dừng với lỗi \texttt{MissingFieldDataError} nêu tên trường, thay vì tự điền giá trị mặc định. Tệp cấu hình hiện trường đo thực mang nhãn \texttt{MEASURED} và hiện chưa có giá trị đo; tệp cấu hình đại diện mang nhãn \texttt{LITERATURE\_PROXY}, trong đó mỗi giá trị đi kèm nguồn và mức bằng chứng. Chương trình không bao giờ tự chuyển từ tệp đo thực sang tệp đại diện.

\textbf{Gói \texttt{net}.} Mạng được dựng bằng hai lượt \texttt{netconvert}: lượt đầu cắt theo bán kính quanh nút và giữ các cạnh cho phép ô tô và xe máy, lượt sau chỉ giữ bốn nhánh và đơn giản hóa hình học. Mô-đun áp vá sinh các tệp plain XML cho bề rộng làn, kết nối rẽ và chương trình đèn từ cấu hình hiện trường, trong đó kết nối rẽ trái trên nhánh cấm rẽ được thêm lại với quyền chỉ dành cho vClass \texttt{custom1}; liên kết rẽ phải của làn rẽ phải liên tục mang trạng thái \texttt{g} ở mọi pha; mỗi nhánh được tách thành đoạn phía trước và đoạn cọc 30~m cấm chuyển làn giữa lane 0 và lane 1; nhánh Lê Văn Lương phía tây có thêm cạnh vỉa hè chỉ dành cho \texttt{custom1} với hai liên kết đèn riêng. Việc tách cạnh làm vạch dừng dịch đi một khoảng nhỏ (từ $-0{,}70$ đến $+1{,}61$~m tùy nhánh), được ghi vào một tệp để các bảng quan sát dùng đúng vị trí; chỉ số 24 liên kết đèn của lòng đường không đổi. Toàn bộ quá trình được gọi qua một lệnh dòng lệnh duy nhất.

\textbf{Gói \texttt{rules}.} Mô-đun lược đồ định nghĩa tham số hợp lệ của từng loại luật TraCI và các vClass mà mỗi luật được áp dụng. Hàm kiểm tra hợp lệ trả về toàn bộ lỗi trong một danh sách, mỗi lỗi mang một trong sáu mã \texttt{schema}, \texttt{unknown\_attr}, \texttt{out\_of\_bounds}, \texttt{share\_sum}, \texttt{rule\_params}, \texttt{logic}; nội dung lỗi nêu rõ khoảng giá trị hoặc tên được phép để có thể đưa nguyên văn vào phản hồi cho LLM, và một nhật ký hợp lệ ghi lại kết quả kiểm tra. Hàm biên dịch gọi bộ kiểm tra trước; nếu đặc tả không hợp lệ, nó báo lỗi \texttt{InvalidSpecError} kèm danh sách lỗi và không ghi tệp nào. Với đặc tả hợp lệ, nó ghi \texttt{vtypes.add.xml} gồm các phân phối \texttt{car}, \texttt{moto}, \texttt{moto\_violator} và (nếu có) \texttt{car\_violator} theo thứ tự cố định, mỗi nguyên mẫu là một vType có mã trùng mã nguyên mẫu và xác suất bằng tỷ trọng; nguyên mẫu \texttt{custom1} (\texttt{custom2}) được gán trước các giá trị mặc định của xe máy (ô tô) rồi mới ghi đè bởi đặc tả, vì các vClass tùy biến không có mặc định này. Mọi tham số, kể cả khi bằng giá trị mặc định, đều được ghi tường minh. Tệp sinh ra khai báo lược đồ XSD của SUMO, nhờ đó SUMO kiểm tra tên và giá trị thuộc tính mỗi khi nạp; thử nghiệm cho thấy nếu không khai báo, thuộc tính sai tên bị bỏ qua mà không có cảnh báo. Tệp \texttt{traci\_rules.json} chỉ được ghi khi đặc tả có luật, và mỗi luật mang danh sách mã vType thành viên mà nó áp dụng, vì xe lấy từ một phân phối báo mã của vType thành viên chứ không phải mã của phân phối.

\textbf{Gói \texttt{runtime}.} Mô-đun trạng thái đọc thông tin nút giao một lần khi bắt đầu, sau đó mỗi bước dùng một đăng ký ngữ cảnh quanh nút với bán kính 220~m và một đăng ký trạng thái đèn để dựng ảnh chụp trạng thái chung cho mọi luật. Lớp \texttt{Controller} tạo một đối tượng luật cho mỗi mục trong \texttt{traci\_rules.json} thông qua một bảng đăng ký (\texttt{RULE\_REGISTRY}) ánh xạ loại luật sang lớp hiện thực. Mỗi luật có bộ sinh số ngẫu nhiên riêng, khởi tạo từ seed của lần chạy, chỉ số và loại luật; xe được duyệt theo thứ tự mã, nên hai backend \texttt{libsumo} và \texttt{traci} cho cùng chuỗi rút ngẫu nhiên và cùng nhật ký can thiệp. Mọi lệnh gửi tới SUMO đi qua lớp bao \texttt{Commands}, đảm bảo mỗi xe chỉ chịu một can thiệp tại một thời điểm, chỉ xóa đúng bit được nêu tên của speed mode và kiểm tra lại giá trị SUMO báo về, đồng thời lưu các khoảng thời gian xe chịu can thiệp. Một khoảng như vậy chỉ đóng khi xe đã đi qua nút (xuất hiện trên làn không thuộc nút và không phải làn vào) hoặc rời vùng xét, để va chạm trong nút của xe vượt đèn đỏ vẫn được tính là liên quan đến can thiệp. Bộ điều khiển cung cấp bảng tóm tắt số can thiệp đã thử và đã thực thi theo từng luật, thời gian xử lý của nó, và hàm cho biết một xe có đang trong khoảng can thiệp tại một thời điểm hay không.

\textbf{Gói \texttt{sim}.} Mô-đun cấu hình ghi tệp \texttt{run.sumocfg} bằng chính SUMO (tùy chọn lưu cấu hình), nên tùy chọn sai tên bị SUMO từ chối; mô-đun này luôn bật kiểm tra va chạm trong nút (\texttt{collision.check-junctions}) và gộp cảnh báo cùng loại. Trình chạy cung cấp hàm chạy một lần, hàm chạy nhiều lần song song và giao diện dòng lệnh. Tệp \texttt{manifest.json} của mỗi lần chạy ghi: mã lần chạy, mã đặc tả, ngày, seed; mã commit và trạng thái thay đổi chưa commit; phiên bản SUMO và backend; giá trị băm của đặc tả, mạng, tệp tuyến và cấu hình hiện trường; nhãn nguồn cấu hình hiện trường; cấu hình thí nghiệm hiệu lực và tùy chọn SUMO; thời gian chạy; số teleport (tổng và không do va chạm); số va chạm (tổng, theo loại, trên làn và trong nút, liên quan can thiệp TraCI); tóm tắt can thiệp và thời gian xử lý của bộ điều khiển; số xe chưa chèn được vào mạng (cuối lần chạy và lớn nhất); số cảnh báo và lỗi; và cờ \texttt{valid\_strict}, chỉ đúng khi không có teleport và không có va chạm nào. Một hàm riêng đánh giá tính đủ điều kiện của lần chạy theo các ngưỡng trong cấu hình thí nghiệm. Mô-đun đầu ra đọc đầu ra SUMO (cả \texttt{.xml} và \texttt{.xml.gz}) trong một lượt duyệt dữ liệu quỹ đạo và ghi mười bảng theo lược đồ quan sát: số đếm rẽ, sự kiện đèn, lượt qua vạch dừng, ảnh chụp hàng chờ, vị trí ngang, hàng chờ giữa nút, lượt đi qua nút, chuỗi thời gian hàng chờ (bước 2~s), sự kiện tìm khoảng trống và độ trễ xe rẽ phải. Ba bảng cuối chỉ được suy từ quỹ đạo, detector và tuyến, không đọc nhật ký can thiệp, để cùng một định nghĩa áp dụng được cho dữ liệu video.

\textbf{Gói \texttt{proxy} và \texttt{metrics}.} Bộ sinh dữ liệu đại diện tạo, với mỗi ngày, mười bảng CSV cùng một tệp manifest và tệp tuyến tương ứng; nó không đọc đầu ra SUMO, và ba bảng của nhóm F2 được sinh bằng một luồng số ngẫu nhiên riêng để các bảng cũ không đổi. Gói chỉ số gồm một mô-đun dùng chung (gán chu kỳ đèn theo mọi mốc bắt đầu xanh, xác định xe rẽ phải liên tục để loại khỏi các chỉ số vi phạm), một mô-đun cho mỗi chỉ số trong Bảng~\ref{table:3.3}, mô-đun báo cáo xác suất chọn không gian thay thế theo độ dài hàng chờ, và mô-đun khoảng cách. Hàm \texttt{compare} của mô-đun khoảng cách trả về một báo cáo gồm $S$, $S_{\text{sourced}}$, $S_{\text{assumed}}$, khoảng cách và tỷ số $R$ theo từng khóa, cùng các dòng tóm tắt không chứa tên ngày để dùng trực tiếp trong phản hồi cho LLM.

\textbf{Mô-đun \texttt{evaluate}.} Lớp \texttt{RunBudget} giữ ngân sách của một chiến dịch trong tệp \texttt{budget.json}, đọc lại từ đĩa và trừ dưới khóa tệp. Hàm \texttt{evaluate\_spec} kiểm tra ngày, kiểm tra hợp lệ và biên dịch đặc tả vào thư mục của chiến dịch, trừ ngân sách, chạy song song các lần chạy, áp cổng G0 và tiêu chí đủ điều kiện, rồi gọi hàm khoảng cách trên mọi lần chạy. Đặc tả không hợp lệ không tốn lần chạy; mã đặc tả trùng trong cùng chiến dịch bị từ chối; ngày đánh giá bị chặn trừ khi được gọi tường minh từ bước đánh giá cuối. Hàm \texttt{select\_best} cài đặt quy tắc chọn đặc tả tốt nhất ở Mục~\ref{section:3.8}.

\textbf{Gói \texttt{llm}.} Lớp \texttt{LLMClient} có hai backend: \texttt{openai} gọi máy chủ tương thích OpenAI (yêu cầu khai báo nhiệt độ và mức lượng tử hóa), và \texttt{file} ghi yêu cầu ra tệp rồi chờ tệp phản hồi, dùng cho mô hình thay thế. Với backend \texttt{file}, vòng lặp gửi mọi yêu cầu của một vòng, dừng với mã thoát riêng khi còn thiếu phản hồi, và được chạy tiếp sau khi các phản hồi đã có. Khi phát lại nhật ký, trình khách kiểm tra định danh LLM để một câu trả lời của mô hình thay thế không thể bị ghi nhận là của mô hình khác. Hàm \texttt{run\_loop} lưu trạng thái của chiến dịch (chế độ, đặc tả cơ sở, ngày, seed, ngân sách, thời lượng mô phỏng, cấu hình hiện trường, mức trần va chạm B0 và định danh LLM) và từ chối chạy tiếp nếu các giá trị này thay đổi; mỗi vòng ghi một tệp \texttt{iter\_<i>.json} với tỷ lệ hợp lệ và đặc tả tốt nhất đến lúc đó.

\textbf{Gói \texttt{baselines}.} Hàm \texttt{search\_space} dựng không gian Optuna từ cấu trúc cho trước (Mục~\ref{section:3.8}); nghiên cứu Optuna được lưu trong một tệp nhật ký nên có thể chạy tiếp mà không trừ lại ngân sách, và từ chối chạy tiếp nếu ngày, seed hoặc cấu trúc khác. Đặc tả tốt nhất được chọn bằng cùng hàm \texttt{select\_best} như vòng lặp LLM và mang \texttt{provenance.model} là \texttt{optuna-tpe}.

\section{Thiết kế chi tiết}

\subsection{Thiết kế giao diện}

\subsection{Thiết kế lớp}

Mỗi luật TraCI là một lớp con của lớp cơ sở \texttt{Rule}. Lớp cơ sở giữ tham số của luật, tập mã vType áp dụng, thông tin nút giao, lớp bao lệnh và bộ sinh số ngẫu nhiên, cùng một từ điển các can thiệp đang diễn ra (mỗi phần tử ghi mã xe, hành động, thời điểm bắt đầu, speed mode gốc nếu đã bị thay đổi và cờ cho biết tốc độ đang bị ghi đè). Phương thức trừu tượng \texttt{on\_step} nhận ảnh chụp trạng thái của bước hiện tại và trả về danh sách các can thiệp đã có kết quả trong bước đó. Lớp cơ sở cung cấp các thao tác dùng chung: kiểm tra xe có thuộc phạm vi luật và đang rảnh hay không, bắt đầu can thiệp, hoàn tác ghi đè (trả tốc độ về cho SUMO quyết định và khôi phục speed mode gốc), kết thúc can thiệp với kết quả và lý do, và đóng mọi can thiệp còn dang dở khi lần chạy kết thúc với lý do \texttt{run\_ended}. Khi SUMO báo speed mode khác với giá trị bộ điều khiển vừa đặt, lớp bao lệnh báo lỗi \texttt{ModeNotRestoredError} thay vì bỏ qua.

Kết quả của một can thiệp là một bản ghi bất biến gồm bảy trường: thời điểm phát lệnh, mã xe, loại luật, hành động, cờ thực thi (phản ứng quan sát được của SUMO, định nghĩa riêng cho từng luật), lý do không thực thi (rỗng nếu đã thực thi) và thời điểm kết quả được xác định. Mỗi bản ghi là một dòng của \texttt{interventions.csv}.

\subsection{Thiết kế cơ sở dữ liệu}

\section{Xây dựng ứng dụng}

\subsection{Thư viện và công cụ sử dụng}

Bảng~\ref{table:4.3} liệt kê các công cụ và thư viện. Với các thư viện Python, cột phiên bản ghi ràng buộc phiên bản tối thiểu khai báo trong cấu hình gói của dự án; phiên bản thực tế đã cài trong môi trường chạy chưa được ghi lại \textbf{[TODO: ghi phiên bản đã cài từ môi trường ảo]}.

\begin{table}[H]
\centering
\small
\begin{tabular}{|p{5.0cm}|p{4.2cm}|p{3.6cm}|}
\hline
\textbf{Mục đích} & \textbf{Công cụ} & \textbf{Phiên bản} \\ \hline
Mô phỏng vi mô, dựng mạng, sinh tuyến & SUMO (\texttt{sumo}, \texttt{sumo-gui}, \texttt{netconvert}, \texttt{routeSampler})~\cite{lopez2018sumo} & 1.27.0 \\ \hline
Điều khiển SUMO từ Python, chạy hàng loạt & \texttt{libsumo}, \texttt{sumo-data} & 1.27.0 \\ \hline
Điều khiển SUMO qua giao diện đồ họa, đọc mạng & \texttt{traci}, \texttt{sumolib} (lấy từ thư mục cài đặt SUMO) & Cùng phiên bản SUMO (1.27.0) \\ \hline
Ngôn ngữ lập trình & Python (môi trường ảo) & 3.12 \\ \hline
Đọc cấu hình & \texttt{pyyaml} & $\geq$ 6.0 \\ \hline
Tính toán số, xử lý bảng, thống kê & \texttt{numpy}, \texttt{pandas}, \texttt{scipy} & $\geq$ 1.26; $\geq$ 2.1; $\geq$ 1.11 \\ \hline
Đọc và ghi XML & \texttt{lxml} & $\geq$ 5.0 \\ \hline
Vẽ hình & \texttt{matplotlib} & $\geq$ 3.8 \\ \hline
Kiểm tra đặc tả theo lược đồ JSON & \texttt{jsonschema} & $\geq$ 4.20 \\ \hline
Chuyển tọa độ kinh vĩ độ sang tọa độ mạng & \texttt{pyproj} & $\geq$ 3.6 \\ \hline
Kiểm thử tự động & \texttt{pytest} & $\geq$ 8.0 \\ \hline
Gọi LLM qua giao diện tương thích OpenAI & \texttt{openai} & $\geq$ 1.30 \\ \hline
Tối ưu tham số B2, H~\cite{akiba2019optuna} & \texttt{optuna} & $\geq$ 3.5 \\ \hline
Chạy LLM cục bộ (chỉ trên máy Windows; chưa dùng cho chiến dịch nào) & Ollama, mô hình \texttt{qwen3.5:9b} & Ollama 0.40.10 \\ \hline
Mô hình thay thế cho LLM cục bộ trong các chiến dịch P1, P2 hiện tại & Claude Sonnet (qua backend tệp) & \textbf{[TODO: phiên bản mô hình]} \\ \hline
Quản lý phiên bản mã & Git & \textbf{[TODO]} \\ \hline
\end{tabular}
\caption{Thư viện và công cụ sử dụng}
\label{table:4.3}
\end{table}

Thư viện \texttt{traci} và \texttt{sumolib} không được cài từ kho PyPI mà được nạp từ thư mục công cụ của bản cài SUMO, để luôn khớp với tệp thực thi SUMO. Gói \texttt{libsumo} và \texttt{sumo-data} được cài không kèm phụ thuộc, vì phụ thuộc khai báo của chúng kéo theo \texttt{traci} và \texttt{sumolib} phiên bản khác từ PyPI và che mất bản đi kèm SUMO. Hai backend được kiểm tra cho kết quả giống nhau bằng một kiểm thử riêng (Mục~\ref{section:testing}). Mô hình LLM đích theo kế hoạch là một mô hình lớn hơn; mô hình \texttt{qwen3.5:9b} chỉ là lựa chọn tạm thời cho phần cứng của máy Windows, và mức lượng tử hóa của nó chưa được chốt nên cấu hình này chưa được bật. Các chiến dịch hiện tại chạy trên máy Linux, nơi chưa có runtime LLM cục bộ, nên cấu hình LLM đang hoạt động là mô hình thay thế với 4 mẫu mỗi vòng và tối đa 3 lần thử \textbf{[TODO: chốt mô hình LLM cục bộ]}.

\subsection{Kết quả đạt được}

Công việc được chia thành 15 hạng mục. Bảng~\ref{table:4.4} tóm tắt sản phẩm và trạng thái của từng hạng mục tại thời điểm viết.

\begin{table}[H]
\centering
\small
\begin{tabular}{|p{0.8cm}|p{9.6cm}|p{2.6cm}|}
\hline
\textbf{\#} & \textbf{Sản phẩm} & \textbf{Trạng thái} \\ \hline
1 & Đọc cấu hình; tệp cấu hình hiện trường đại diện có nguồn và mức bằng chứng cho từng giá trị & Hoàn thành \\ \hline
2 & Mạng mô phỏng sinh bằng script từ OSM và tệp vá; bảng cạnh; sơ đồ mạng & Hoàn thành \\ \hline
3 & vType B0, sinh demand, cấu hình SUMO và các đầu ra & Hoàn thành \\ \hline
4 & Trình chạy, manifest, chuyển đầu ra sang lược đồ quan sát & Hoàn thành \\ \hline
4b & Sửa bề rộng làn ra; tắt mặc định thiết bị đo chỉ số an toàn thay thế (SSM) của SUMO, gộp cảnh báo; chạy lại B0 trên dữ liệu đại diện & Hoàn thành \\ \hline
5 & Thí nghiệm khả thi về xe chờ giữa nút; quyết định dùng cơ chế sẵn có của SUMO & Hoàn thành \\ \hline
6 & Bảng giá trị mục tiêu và bộ sinh dữ liệu đại diện PROXY\_D1--PROXY\_D5 & Hoàn thành \\ \hline
7 & Miền tham số, lược đồ đặc tả luật, bộ kiểm tra hợp lệ, đặc tả B1 & Hoàn thành \\ \hline
8 & Bộ biên dịch đặc tả & Hoàn thành \\ \hline
9 & Bộ điều khiển TraCI và năm luật có điều kiện & Hoàn thành \\ \hline
10 & Các chỉ số M1--M8, G0 và hàm khoảng cách & Hoàn thành \\ \hline
F2 & Mạng F2 (làn rẽ phải liên tục, \texttt{custom2}), thí nghiệm khả thi Phase 3A--3D, đoạn cọc, vỉa hè, luật \texttt{queue\_space\_seeking}, chỉ số M9--M10, đóng băng ngưỡng va chạm trên mạng F2 & Hoàn thành (đi ngược chiều được hoãn) \\ \hline
11 & Khung đánh giá chung và bộ đếm ngân sách số lần chạy & Hoàn thành \\ \hline
12 & Gọi LLM, prompt, phản hồi & Hoàn thành (chưa thử với LLM cục bộ) \\ \hline
13 & Vòng lặp LLM P1 (\texttt{params\_only}) và P2 (\texttt{full}) & Mã hoàn thành; chiến dịch P1, P2 tạm dừng sau 2/3 vòng, 24/36 lần chạy có kết quả (mô hình thay thế) \\ \hline
14 & Phương án đối sánh B2 và H bằng Optuna & Mã hoàn thành; B2 36/36 lần chạy; H chưa chạy \\ \hline
15 & Đánh giá cuối và báo cáo & Dự kiến \\ \hline
\end{tabular}
\caption{Các hạng mục và trạng thái}
\label{table:4.4}
\end{table}

Thống kê kích thước mã nguồn chưa được tổng hợp \textbf{[TODO: số tệp và số dòng mã của từng gói]}.

\subsection{Minh họa các chức năng chính}

Quá trình dựng mạng sinh ra một sơ đồ mạng dạng SVG, và thí nghiệm về xe chờ giữa nút sinh hai hình về số xe trong vùng chờ và vị trí ngang của các xe chờ \textbf{[TODO: chuyển sơ đồ mạng và hai hình của thí nghiệm vùng chờ sang PNG, chèn vào mục này kèm chú thích]}. Ảnh chụp \texttt{sumo-gui} khi xe chờ giữa nút cũng đã được ghi nhận trong quá trình điều tra va chạm \textbf{[TODO: bổ sung ảnh chụp sumo-gui]}.

\section{Kiểm thử}
\label{section:testing}

Kiểm thử tự động dùng \texttt{pytest}. Phần lớn kiểm thử chạy trên dữ liệu nhỏ dựng tay hoặc trên các mạng chữ thập tí hon, với giá trị kỳ vọng tính độc lập với mã được kiểm thử; một số kiểm thử chạy SUMO thật trong thời gian ngắn. Bảng~\ref{table:4.5} tóm tắt các bộ kiểm thử.

\begin{table}[H]
\centering
\small
\begin{tabular}{|p{3.6cm}|p{9.4cm}|}
\hline
\textbf{Bộ kiểm thử} & \textbf{Nội dung kiểm tra} \\ \hline
\texttt{test\_env} & \texttt{traci} và \texttt{sumolib} được nạp từ bản cài SUMO và khớp phiên bản tệp thực thi \\ \hline
\texttt{test\_config} & Trường bắt buộc rỗng gây lỗi nêu tên trường; cấu hình đo thực không bao giờ tự lấy giá trị đại diện \\ \hline
\texttt{test\_net} & Dựng mạng bằng một lệnh; tệp vá, đèn, bảng cạnh; dựng lại cho kết quả giống từng byte; làn nội bộ có bề rộng bằng làn ra tương ứng \\ \hline
\texttt{test\_demand} & vType B0, tuyến ứng viên, sinh tuyến từ số đếm rẽ và kiểm tra dữ liệu số đếm \\ \hline
\texttt{test\_sumocfg} & Tùy chọn SUMO lấy từ cấu hình thí nghiệm, tệp đầu ra, chế độ SSM, gộp cảnh báo, chạy thật ngắn \\ \hline
\texttt{test\_runner} & Trình chạy, các trường của manifest, chuyển đổi sang lược đồ quan sát, tiêu chí đủ điều kiện \\ \hline
\texttt{test\_outputs\_\allowbreak synthetic} & Chuyển đổi lược đồ trên đầu ra SUMO viết tay cho mạng chữ thập nhỏ, kết quả kỳ vọng tính tay \\ \hline
\texttt{test\_proxy} & Tính tất định của bộ sinh dữ liệu đại diện, lược đồ, tính nhất quán nội tại và độ khớp với giá trị mục tiêu \\ \hline
\texttt{test\_validator} & Miền tham số, tham số luật, các mã lỗi và nội dung thông báo lỗi của bộ kiểm tra hợp lệ \\ \hline
\texttt{test\_compiler} & Biên dịch B1 khớp từng byte với bản chụp; mọi thuộc tính nằm trong miền; SUMO nạp được với kiểm tra XSD; đặc tả B0 tương đương tệp vType viết tay; mã vType xe báo về khi lấy từ phân phối \\ \hline
\texttt{test\_runtime} & Mỗi luật một kịch bản trên mạng nhỏ; lệnh bị SUMO bỏ qua được ghi là không thực thi; khôi phục speed mode; tích hợp với trình chạy; giữ xe giữa nút không gây va chạm trong nút; cấm \texttt{ignoreIDs} trong mã nguồn \\ \hline
\texttt{test\_backends} & \texttt{libsumo} và \texttt{traci} cho quỹ đạo giống hệt nhau với cùng kịch bản và seed \\ \hline
\texttt{test\_metrics} & Các chỉ số M1--M10, G0 và hàm khoảng cách trên bảng dựng tay có đáp án; khoảng cách khi một phía rỗng không bao giờ nhỏ hơn khoảng cách khi sai \\ \hline
\texttt{test\_evaluate} & Ngân sách theo chiến dịch không vượt khi chạy lại; đặc tả hỏng không tốn lần chạy; mã đặc tả trùng bị từ chối; chặn ngày đánh giá; quy tắc chọn đặc tả tốt nhất \\ \hline
\texttt{test\_llm} & Lược đồ cho LLM, khóa cấu trúc, phát lại nhật ký có kiểm định danh LLM, chặn rò rỉ dữ liệu ngày đánh giá trong phản hồi \\ \hline
\texttt{test\_loop} & Vòng lặp chạy tiếp đúng sau khi bị ngắt (giữa vòng, ở đặc tả cuối còn đủ ngân sách) mà không trừ lại ngân sách; mẫu không hợp lệ không tốn lần chạy; khóa cấu trúc của P1; chạy thêm vòng khi còn ngân sách; số lần chạy đã dùng khớp bộ đếm ngân sách \\ \hline
\texttt{test\_optuna} & Không gian tìm kiếm theo cấu trúc; mọi thử nghiệm ngẫu nhiên đều hợp lệ; chạy tiếp nghiên cứu không trừ lại ngân sách \\ \hline
\end{tabular}
\caption{Các bộ kiểm thử tự động}
\label{table:4.5}
\end{table}

Lần chạy toàn bộ kiểm thử gần nhất được ghi nhận, sau khi hoàn thành vòng lặp LLM, cho kết quả 367 kiểm thử đạt và 3 kiểm thử được đánh dấu là thất bại dự kiến (xfail). Cả ba kiểm thử xfail so sánh một tỷ lệ trên năm ngày đại diện với giá trị mục tiêu, với ngưỡng chấp nhận $\pm$15\%: tỷ lệ vượt đỏ của xe máy nhánh Lê Văn Lương phía tây lệch $+17{,}6\%$, tỷ lệ xuất phát sớm của ô tô lệch $-19\%$ và tỷ lệ vượt đỏ của ô tô lệch $-45\%$. Các độ lệch này là sai số ngẫu nhiên của bộ năm ngày (ô tô có rất ít sự kiện; với xe máy, trung bình của 200 bộ sinh lại chỉ lệch $+0{,}1\%$), và bộ ngày không được sinh lại theo kết quả.

Hai kiểm thử có vai trò đặc biệt đối với độ tin cậy của kết quả. Kiểm thử đối chiếu backend đảm bảo việc dùng \texttt{libsumo} để tăng tốc không làm thay đổi kết quả so với \texttt{traci}; ở cấp bộ điều khiển, hai backend cũng cho cùng nhật ký can thiệp. Kiểm thử giữ xe giữa nút chạy trên mạng chữ thập nhỏ có dòng xe cắt ngang và bật kiểm tra va chạm trong nút: giữ xe thuần bằng cách đặt tốc độ bằng 0 không gây va chạm trong nút, còn một đối chứng âm cố ý cho xe cắt ngang bỏ qua xe đang bị giữ thì gây va chạm, nên kiểm thử phát hiện được nếu cách làm sai này quay lại. Script thí nghiệm vùng chờ giữa nút (Mục~\ref{section:box_capacity}) không có kiểm thử tự động.

\section{Thực nghiệm}
\label{section:experiments}

Mục này trình bày các thực nghiệm đã tiến hành. Tất cả đều chạy trên mạng đại diện (cấu hình hiện trường lấy từ tài liệu và giả định). Các Mục~\ref{section:box_capacity}--\ref{subsection:b1_check} được thực hiện trên mạng trước khi bổ sung làn rẽ phải liên tục (gọi là mạng trước F2); các Mục~\ref{subsection:phase3}--\ref{subsection:interim} dùng mạng cuối (mạng F2). Số liệu của hai mạng không so sánh trực tiếp được với nhau. Với mỗi thực nghiệm, phần quan sát chỉ nêu số liệu, phần diễn giải chỉ dựa trên các số liệu đó, còn các giả thuyết cần kiểm chứng được nêu riêng. Trong mục này, \emph{va chạm/h} là số va chạm SUMO ghi nhận trong một lần chạy quy về một giờ mô phỏng; \emph{teleport kẹt} là số lần SUMO đưa một xe ra khỏi mạng vì xe bị kẹt quá thời gian cho phép (mặc định 300~s) hoặc kẹt khi nhường đường, không tính teleport do va chạm.

\subsection{Thí nghiệm khả thi về xe chờ giữa nút}
\label{section:box_capacity}

\textbf{Câu hỏi và thiết lập.} Thí nghiệm trả lời câu hỏi: với mô hình làn con SL2015, SUMO 1.27.0 có tự mô phỏng được cảnh xe máy rẽ trái trái phép dừng chờ giữa nút hay không, và chứa được bao nhiêu xe. Nếu sức chứa tối đa của cơ chế sẵn có không nhỏ hơn trung bình mục tiêu của số xe chờ giữa nút (chỉ số M7a, Bảng~\ref{table:3.3}), bằng 6 xe mỗi chu kỳ, thì dùng cơ chế sẵn có; nếu không, phải sửa mạng bằng một cạnh và tín hiệu riêng cho vùng chờ. Thí nghiệm dùng mạng đại diện với chu kỳ đèn 120~s (xanh, vàng, đỏ toàn phần lần lượt 64/4/2~s cho Lê Văn Lương và 44/4/2~s cho Hoàng Đạo Thúy) và vType B0. Demand được tổng hợp riêng cho thí nghiệm, không phải dữ liệu đại diện: dòng đi thẳng đối diện dày gồm 4000 xe máy/h và 700 ô tô/h (không còn khoảng trống trong pha xanh), hoặc vừa phải gồm 1300 xe máy/h và 200 ô tô/h (gần lưu lượng đi thẳng của nhánh đối diện trong các ngày đại diện); mỗi nhánh Hoàng Đạo Thúy có 1200 xe máy/h và 250 ô tô/h đi thẳng. Năm kịch bản khác nhau ở dòng xe vi phạm (Bảng~\ref{table:4.box}): \texttt{sat} với 1500 xe/h làm bão hòa vùng chờ; \texttt{target} với 225 xe/h, tức 7,5 xe vào mỗi chu kỳ theo trung bình mục tiêu; \texttt{sat\_tight} với 1500 xe/h và xe vi phạm có \texttt{minGap} 1,0~m, \texttt{minGapLat} 0,3~m; \texttt{sat\_mg18} với 1500 xe/h và \texttt{minGap} 1,8~m; \texttt{target\_modopp} giống \texttt{target} nhưng với dòng đối diện vừa phải. Mỗi kịch bản chạy 5 seed (1--5), mỗi lần 900~s (7 lần kết thúc pha xanh đối diện), tổng 25 lần chạy. Một bộ quan sát TraCI không can thiệp ghi lại tại mỗi bước 0,1~s vị trí, vị trí ngang, tốc độ của mọi xe trong vùng chờ, trạng thái đèn, teleport và va chạm. Xe được coi là đang dừng khi tốc độ dưới 5~km/h; \emph{sức chứa} là số xe đang dừng lớn nhất trong vùng chờ khi dòng đối diện còn xanh; \emph{thoát sớm} là xe rời vùng chờ qua khe trống trước khi pha xanh đối diện kết thúc, tỷ lệ tính trên số xe vào vùng chờ. Thí nghiệm này chạy trước khi bật kiểm tra va chạm giữa các làn nội bộ cắt nhau.

\begin{table}[H]
\centering
\small\setlength{\tabcolsep}{3pt}
\begin{tabular}{|p{2.3cm}|p{1.5cm}|p{1.3cm}|p{1.2cm}|p{0.9cm}|p{0.9cm}|p{2.4cm}|p{1.6cm}|}
\hline
\textbf{Kịch bản} & \textbf{Đối diện (xe/h)} & \textbf{minGap (m)} & \textbf{Chờ TB} & \textbf{p90} & \textbf{Max} & \textbf{Thoát sớm mỗi chu kỳ (tỷ lệ)} & \textbf{Khoảng xả (s)} \\ \hline
\texttt{sat} & 4700 & 2,5 & 8,00 & 8 & 8 & 0 (0\%) & 0,52 \\ \hline
\texttt{target} & 4700 & 2,5 & 6,43 & 8 & 8 & 0 (0\%) & 0,68 \\ \hline
\texttt{sat\_tight} & 4700 & 1,0 & 12,00 & 12 & 12 & 0 (0\%) & 0,51 \\ \hline
\texttt{target\_\allowbreak modopp} & 1500 & 2,5 & 2,60 & 6 & 7 & 4,17 (62\%) & 0,88 \\ \hline
\texttt{sat\_mg18} & 4700 & 1,8 & 10,00 & 10 & 10 & 0 (0\%) & 0,52 \\ \hline
\end{tabular}
\caption{Xe chờ giữa nút trên vùng chờ rộng 3,2~m (gộp chu kỳ 1--6 của 5 seed, 30 chu kỳ mỗi kịch bản; ``Chờ'' là số xe đang chờ khi pha xanh đối diện kết thúc; ``Khoảng xả'' là khoảng thời gian trung bình giữa hai xe liên tiếp rời vùng chờ)}
\label{table:4.box}
\end{table}

\textbf{Quan sát (vùng chờ 3,2~m).} Trong mọi chu kỳ của mọi lần chạy, xe vi phạm dừng chờ trên làn nội bộ rẽ trái, xe đầu dừng cách điểm chờ nội bộ 0,10~m; khi làn này đầy, các xe tiếp theo dừng tại vạch dừng. Với vType B0, sức chứa đo được là 8 xe ở cả 5 seed, xếp thành hai cột ngang ở vị trí ngang khoảng $-1{,}15$~m và từ 0 đến $+0{,}65$~m. Giảm \texttt{minGap} của xe vi phạm xuống 1,8~m và 1,0~m cho sức chứa 10 và 12 xe. Ở kịch bản \texttt{target}, vùng chờ đầy trong 43\% số chu kỳ. Xe chờ đầu tiên rời vùng chờ trung bình 4,7~s sau khi pha xanh đối diện kết thúc; vùng chờ trống sau 8,2~s với 8 xe và 10,4~s với 12 xe. Khi pha xanh Hoàng Đạo Thúy bắt đầu, vẫn còn 3,4 (\texttt{target}) đến 4,6 (\texttt{sat}) xe trong vùng chờ, làm lượt qua vạch đầu tiên của làn 2 nhánh Hoàng Đạo Thúy phía bắc trễ 0,4--1,1~s với B0 và 1,6--2,7~s với 12 xe. Không có kẹt cứng, không có teleport không do va chạm trong cả 25 lần chạy, thời gian lưu lại vùng chờ lớn nhất là 70,4~s. Tổng cộng có 334 va chạm trong 25 lần chạy, chỉ một va chạm có xe vi phạm tham gia; còn lại là giữa các xe đi thẳng cùng dòng.

\textbf{Diễn giải.} Cơ chế điểm chờ nội bộ sẵn có của SUMO tái hiện được hành vi vào nút, dừng chờ, chặn xe sau khi vùng chờ đầy và xả sau khi dòng đối diện dừng, mà không cần tùy chọn hay sửa mạng. Sức chứa do hình học quyết định: số hàng theo công thức $\lfloor (18{,}79 + \text{minGap}) / (2{,}2 + \text{minGap}) \rfloor$, trong đó 18,79~m là vị trí đầu xe dẫn đầu, dự đoán 4, 5 và 6 hàng tại \texttt{minGap} 2,5; 1,8 và 1,0~m, tức 8, 10 và 12 xe với hai cột, đúng như đo được. Vì sức chứa 8 không nhỏ hơn trung bình mục tiêu 6, đồ án chọn dùng cơ chế sẵn có. Tuy nhiên, sức chứa 8 nhỏ hơn phân vị 90 của mục tiêu (10 xe), nên phân phối số xe chờ bị cắt ở 8 với vType B0; đây là vấn đề hiệu chỉnh tham số (\texttt{minGap} của nguyên mẫu \texttt{custom1}) chứ không phải vấn đề mạng. Tỷ lệ thoát sớm phụ thuộc mạnh vào dòng đối diện: 0\% với dòng dày và 62\% với dòng gần mức ngày đại diện, trong khi mục tiêu là 20\% (giả định); do đó ở mức lưu lượng đại diện, tham số cần hiệu chỉnh chủ yếu phải giảm thoát sớm. Các va chạm của xe đi thẳng thuộc vấn đề bề rộng làn trình bày ở Mục~\ref{subsection:exit_width}, không liên quan vùng chờ.

\textbf{Kiểm tra lại trên vùng chờ 3,5~m.} Sau khi bề rộng làn ra được sửa thành 3,5~m (Mục~\ref{subsection:exit_width}), làn nội bộ rẽ trái cũng rộng 3,5~m, nên kịch bản \texttt{target} được chạy lại với 5 seed, mỗi lần 900~s, vẫn chưa bật kiểm tra va chạm giữa các làn nội bộ cắt nhau. Kết quả ở Bảng~\ref{table:4.box_r21}.

\begin{table}[H]
\centering
\small
\begin{tabular}{|p{1.4cm}|p{3.4cm}|p{3.6cm}|p{2.3cm}|p{1.8cm}|}
\hline
\textbf{Seed} & \textbf{Số xe dừng lớn nhất trong vùng chờ} & \textbf{Số xe chờ TB khi hết xanh đối diện} & \textbf{Bề rộng dàn ngang (m)} & \textbf{Va chạm} \\ \hline
1 & 13 & 6,86 & 2,60 & 0 \\ \hline
2 & 8 & 6,00 & 2,60 & 0 \\ \hline
3 & 8 & 6,57 & 2,60 & 0 \\ \hline
4 & 7 & 4,86 & 2,50 & 1 \\ \hline
5 & 10 & 6,86 & 2,60 & 0 \\ \hline
\end{tabular}
\caption{Kiểm tra lại kịch bản \texttt{target} trên vùng chờ rộng 3,5~m (vType B0, 900~s mỗi seed)}
\label{table:4.box_r21}
\end{table}

\textbf{Quan sát.} Số xe dừng lớn nhất là 13, 8, 8, 7 và 10 theo seed 1--5, xếp thành ba cột ngang thay vì hai; số va chạm mỗi lần chạy là 0--1, so với 10--24 trên vùng chờ 3,2~m ở cùng kịch bản. Mục tiêu M7a có trung bình 6,0, phân vị 90 bằng 10 và phân vị 95 bằng 11.

\textbf{Diễn giải.} Trên vùng chờ 3,5~m, sức chứa không còn là cận trên thấp hơn mục tiêu: giá trị 13 ở seed 1 đã vượt phân vị 95. Quyết định dùng cơ chế sẵn có vẫn giữ nguyên. Các kết luận về sức chứa và tỷ lệ thoát sớm ở trên đo trên vùng chờ 3,2~m; riêng tỷ lệ thoát sớm chưa được đo lại trên vùng chờ 3,5~m với nhiều seed.

\textbf{Giả thuyết cần kiểm chứng.} Sức chứa phụ thuộc trực tiếp vào hình học vùng chờ (chiều dài làn nội bộ, vị trí điểm chờ, bề rộng làn), vốn đang lấy từ tiêu chuẩn thiết kế chứ không phải đo đạc. Khi có hình học đo thực, thí nghiệm cần chạy lại; nếu sức chứa khi đó nhỏ hơn mục tiêu thì phương án sửa mạng cần được xem xét lại.

\subsection{Ảnh hưởng của bề rộng làn ra đến va chạm}
\label{subsection:exit_width}

\textbf{Thiết lập.} Trong phiên bản mạng đầu tiên, cấu hình đại diện chỉ đặt bề rộng 3,5~m cho làn vào; làn ra giữ giá trị mặc định 3,2~m của \texttt{netconvert}, không có nguồn. Cấu hình B0 được chạy trên dữ liệu SYNTHETIC (chỉ dùng để thử đường ống) và trên PROXY\_D1, mỗi trường hợp 3 seed, khi chưa bật kiểm tra va chạm giữa các làn nội bộ cắt nhau. Lần chạy PROXY\_D1 dùng trực tiếp tệp thực thi SUMO, không ghi dữ liệu quỹ đạo, SSM và detector, vì trình chạy khi đó dừng giữa chừng do thiết bị SSM ghi quá nhiều cảnh báo.

\textbf{Quan sát.} Số va chạm/h theo seed 1, 2, 3 là 5, 16, 10 trên SYNTHETIC và 62, 52, 46 trên PROXY\_D1. Trong 31 va chạm của các lần chạy SYNTHETIC, 20 xảy ra khi xe đi vào nút, 8 ở cuối làn gần vạch dừng và 3 là va chạm làn con khi chuyển làn cách vạch dừng 16--18~m. Thử nghiệm với làn ra rộng 3,5~m giảm tổng số va chạm từ 31 xuống 2 trên SYNTHETIC và từ 160 xuống 10 trên PROXY\_D1; tuy nhiên trên PROXY\_D1, số teleport kẹt khi đó là 18, 32 và 10 theo seed. Các yếu tố khác được thử từng yếu tố một: độ phân giải làn con 0,2~m cho 40 va chạm, 0,8~m cho 12 va chạm; bật kiểm tra va chạm trong nút làm số va chạm tăng từ 31 lên 86.

\textbf{Diễn giải.} \texttt{netconvert} gán cho làn nội bộ bề rộng của làn ra. Xe máy trong mô hình làn con có thể đi sát vạch làn (vị trí ngang tới $\pm$1,5~m), nên khi đi từ làn vào 3,5~m sang làn nội bộ 3,2~m mà tâm các làn vẫn cách nhau 3,5~m, các xe máy đi song song chồng lấn nhau. Đây là sản phẩm phụ của một giá trị mặc định do chính mô hình để lại, không phải hành vi giao thông. Cách sửa là khai báo danh sách làn ra trong cấu hình hiện trường với bề rộng bằng làn vào (3,5~m, theo tiêu chuẩn thiết kế đường đô thị \textbf{[TODO: trích dẫn TCVN 13592:2022]}), kèm một kiểm thử khẳng định mọi làn nội bộ có bề rộng bằng làn ra tương ứng. Sau khi sửa, cả 56 làn của mạng rộng 3,5~m. Một lần chạy kiểm tra 600~s trên PROXY\_D1, seed 1, tắt SSM cho 0 va chạm và 0 teleport trên mạng mới, so với 8 va chạm và 8 teleport (đều do va chạm) trên mạng cũ.

\textbf{Giả thuyết cần kiểm chứng.} Kết quả này cho thấy số va chạm của mô hình làn con rất nhạy với bề rộng làn nội bộ. Trên mạng có bề rộng đo thực, nếu làn vào và làn ra khác nhau, va chạm kiểu này có thể xuất hiện lại; khi đó cần báo cáo như một giới hạn của mô hình làn con chứ không coi là hành vi thực.

\subsection{Cấu hình cơ sở B0 trên dữ liệu đại diện}
\label{subsection:b0_proxy}

\textbf{Thiết lập.} Sau khi sửa bề rộng làn ra, B0 được chạy trên ba ngày hiệu chỉnh PROXY\_D1--PROXY\_D3, mỗi ngày 3 seed (1, 2, 3), mỗi lần 1800~s (thay vì 3600~s, vì một lần chạy PROXY\_D1 3600~s mất hơn 20 phút), tắt thiết bị SSM. Các lần chạy này thực hiện \emph{trước} khi đồ án quyết định bật kiểm tra va chạm trong nút cho mọi lần chạy, nên số va chạm không bao gồm va chạm giữa các làn nội bộ cắt nhau. Bảng~\ref{table:4.b0} dùng các đại lượng: teleport kẹt (định nghĩa ở đầu Mục~\ref{section:experiments}); số xe chưa chèn được vào mạng khi kết thúc; cờ \texttt{valid\_strict} (không có teleport và va chạm nào); và tỷ lệ ô có GEH $<5$ của cổng G0 (Bảng~\ref{table:3.3}) tính trên 62 ô theo 15 phút $\times$ nhánh $\times$ hướng $\times$ loại xe, cùng tỷ lệ tương ứng khi gộp loại xe (ô theo nhánh $\times$ hướng).

\begin{table}[H]
\centering
\small\setlength{\tabcolsep}{3pt}
\begin{tabular}{|p{1.3cm}|p{0.9cm}|p{1.5cm}|p{1.6cm}|p{1.8cm}|p{1.8cm}|p{1.6cm}|p{1.7cm}|}
\hline
\textbf{Ngày} & \textbf{Seed} & \textbf{Va chạm/h} & \textbf{Teleport kẹt} & \textbf{Xe chưa chèn cuối} & \textbf{valid\_\allowbreak strict} & \textbf{GEH$<$5 (ô)} & \textbf{GEH$<$5 (nhánh $\times$ hướng)} \\ \hline
D1 & 1 & 2 & 14 & 161 & Không & 0,919 & 0,844 \\ \hline
D1 & 2 & 0 & 22 & 0 & Không & 0,952 & 0,875 \\ \hline
D1 & 3 & 2 & 10 & 0 & Không & 1,000 & 0,969 \\ \hline
D2 & 1 & 2 & 22 & 0 & Không & 1,000 & 1,000 \\ \hline
D2 & 2 & 6 & 23 & 0 & Không & 1,000 & 0,969 \\ \hline
D2 & 3 & 0 & 2 & 0 & Không & 1,000 & 1,000 \\ \hline
D3 & 1 & 0 & 0 & 0 & Có & 1,000 & 1,000 \\ \hline
D3 & 2 & 2 & 11 & 0 & Không & 0,984 & 0,906 \\ \hline
D3 & 3 & 0 & 0 & 0 & Có & 1,000 & 1,000 \\ \hline
\end{tabular}
\caption{B0 trên PROXY\_D1--PROXY\_D3 (1800~s, chưa bật kiểm tra va chạm trong nút, tắt SSM)}
\label{table:4.b0}
\end{table}

\textbf{Quan sát.} Số va chạm/h của 9 lần chạy là 2, 0, 2, 2, 6, 0, 0, 2, 0 (trung bình 1,56); với thời lượng 1800~s, mỗi va chạm tương ứng 2 va chạm/h. So với 46--62 va chạm/h trên mạng làn ra 3,2~m, mức va chạm đã giảm mạnh. Tuy nhiên, 7/9 lần chạy có teleport kẹt và chỉ hai lần chạy (PROXY\_D3 seed 1 và 3) đạt \texttt{valid\_strict}; lần chạy PROXY\_D1 seed 1 kết thúc với 161 xe chưa chèn được vào mạng. Mọi teleport kẹt đều là xe máy; chúng xảy ra thành cụm trong cùng giây, lần đầu hầu như luôn ở vạch dừng làn 2 (làn dùng chung cho đi thẳng, rẽ trái và rẽ phải) của nhánh Lê Văn Lương phía tây rồi lan sang các liên kết làn 2 trong nút. PROXY\_D1 có nhu cầu cao hơn PROXY\_D2 và PROXY\_D3 khoảng 20\% (4556, 3765 và 3857 xe trong 1800~s). Tỷ lệ ô GEH $<5$ theo loại xe nằm trong khoảng 0,919--1,000, đều vượt ngưỡng 0,85; ở mức nhánh $\times$ hướng, tỷ lệ này là 0,844--1,000, giá trị thấp nhất thuộc lần chạy còn xe chưa chèn.

\textbf{Thí nghiệm một yếu tố.} Trên PROXY\_D1 với 3 seed, loại bỏ 217 xe rẽ trái trái phép (2,4\% số xe) làm tổng số teleport còn 0, 0 và 3; loại bỏ mọi xe rẽ trái cho 0, 3 và 2 teleport và không còn teleport kẹt; còn giảm \texttt{minGap} xuống 1,5~m làm kết quả xấu hơn (68 và 103 teleport ở seed 1 và 2).

\textbf{Diễn giải.} Nguyên nhân chính gây kẹt của B0 trên dữ liệu đại diện là các xe máy rẽ trái trái phép chờ giữa nút với vType mặc định, khiến làn 2 của nhánh có xe vi phạm bị tắc. Đây là một phát hiện về hành vi của cấu hình mặc định chứ không phải lỗi mạng, nên mạng không được vá thêm; giảm tình trạng kẹt này là một mục tiêu cải thiện mà các bộ luật hiệu chỉnh cần đạt. Tỷ lệ GEH cao chỉ cho thấy nút phục vụ được nhu cầu đại diện: các số ``quan sát'' đại diện và tệp tuyến cùng được sinh từ một bộ sinh, nên G0 ở đây không phải bằng chứng về độ thực tế.

\textbf{Giới hạn về điều kiện chạy.} Từ 9 tỷ lệ va chạm trên, ngưỡng va chạm cho phép từng được đặt bằng $\max(2, \lceil p_{90} \rceil) = 3$ va chạm/h (với $p_{90} = 2{,}8$). Ngưỡng này \emph{đã bị hủy} khi đồ án quyết định bật kiểm tra va chạm trong nút cho mọi lần chạy, vì khi tắt kiểm tra, va chạm giữa xe đang chờ giữa nút và dòng cắt ngang không được đếm; hai ngưỡng mới (va chạm trên làn và trong nút) được tính lại từ B0 với kiểm tra bật trên mạng cuối và trình bày ở Mục~\ref{subsection:r28_f2}. Do đó, các số va chạm trong Bảng~\ref{table:4.b0} không dùng để so sánh với các cấu hình chạy sau quyết định này. Toàn bộ kết quả trong mục chỉ phản ánh hành vi của mô hình trên dữ liệu đại diện và chỉ bao phủ nửa giờ đầu của khung cao điểm.

\subsection{Kiểm tra chức năng bộ luật thủ công B1}
\label{subsection:b1_check}

\textbf{Thiết lập.} Đặc tả B1 (tham số vType chuyển từ Saval và cộng sự~\cite{saval2023mixed}, kèm luật vượt đèn đỏ theo đoàn với $p_0 = p_{\max} = 0{,}5$, $k = 0$ và luật xuất phát sớm với $p = 0$) được chạy một lần 600~s trên PROXY\_D1, seed 1, đã bật kiểm tra va chạm trong nút, trên máy đang chạy song song các tác vụ khác. Lần chạy này nhằm kiểm tra bộ điều khiển hoạt động đúng trong trình chạy, không phải một đánh giá; chưa có lần chạy B0 cùng điều kiện để so sánh.

\textbf{Quan sát.} Khi chỉ dùng vType B1 và khi thêm bộ điều khiển với các luật không bao giờ kích hoạt, hai lần chạy cho cùng kết quả: 0 va chạm trên làn, 7 va chạm trong nút, 11 teleport trong đó 4 không do va chạm; quỹ đạo giống hệt nhau (628 bản ghi chuyến đi, so sánh không có khác biệt). Với đầy đủ B1 (luật vượt đỏ theo đoàn và xuất phát sớm), lần chạy có 0 va chạm trên làn, 13 va chạm trong nút (1 va chạm liên quan xe đang chịu can thiệp) và 13 teleport, không teleport nào không do va chạm. Luật vượt đỏ theo đoàn đưa ra 80 quyết định, 53 được SUMO thực thi, 27 không thực thi vì pha đỏ đã kết thúc trước khi xe qua vạch; luật xuất phát sớm với $p = 0$ không tạo can thiệp nào qua 69 lần rút ngẫu nhiên. Thời gian xử lý của bộ điều khiển là 3--7~s cho mỗi 600~s mô phỏng.

\textbf{Diễn giải.} Bộ điều khiển không làm thay đổi quỹ đạo khi không có luật nào kích hoạt, ghi nhận đúng các lệnh bị SUMO bỏ qua, và có chi phí tính toán nhỏ. Các số va chạm và teleport chỉ đến từ một seed, 600~s, nên không được dùng để kết luận về chất lượng của B1.

\subsection{Thí nghiệm khả thi Phase 3A--3D trên mạng có làn rẽ phải liên tục}
\label{subsection:phase3}

\textbf{Bối cảnh.} Khi xác nhận mọi nhánh có làn rẽ phải liên tục, lane 0 của mỗi nhánh trở thành làn chỉ để rẽ phải (không thêm làn mới), nên mỗi nhánh còn hai làn đi thẳng có đèn. Trước khi cài đặt các hành vi (6)--(8) và xác nhận lại hành vi (10) trên mạng này, bốn thí nghiệm khả thi được thực hiện. Ba thí nghiệm 3A, 3B, 3C chạy trên các kịch bản đồ chơi riêng, nên mọi con số dưới đây là thiết lập thí nghiệm, không phải số đo; thí nghiệm 3D chạy trên mạng đại diện F2. Mọi lần chạy bật kiểm tra va chạm trong nút.

\textbf{Quan sát.}
\begin{itemize}
\item \textbf{3A -- chặn làn rẽ phải (hành vi 6).} Mô hình SL2015 không tự tạo hành vi này: trong 12 lần chạy với hàng chờ tới 232~m, không xe vi phạm nào đi vào lane 0, vì với tuyến đi thẳng lane 0 là ngõ cụt. Hành vi chỉ xuất hiện khi có một điều kiện kích hoạt qua TraCI theo độ dài hàng chờ, kết hợp với việc tách đoạn có cọc thành cạnh riêng cấm chuyển làn (giá trị \texttt{ignoring} của thuộc tính cấm chuyển làn bị \texttt{netconvert} 1.27 bỏ qua mà không báo, nên dùng \texttt{emergency}). Khi đó xe vi phạm đi trên lane 0 và dừng trước đầu đoạn cọc để chờ nhập lại; xe rẽ phải bị trễ 4,2~s với xác suất kích hoạt 0,3 và 8--27~s với xác suất 1, so với 3,4--3,6~s khi không có xe chặn. Không có va chạm trong nút trong 90 lần chạy; va chạm trên làn tăng theo xác suất kích hoạt (0--4 mỗi lần chạy); xe tuân thủ không bị kẹt.
\item \textbf{3B -- đi trên vỉa hè (hành vi 7).} SL2015 không tự dùng làn vỉa hè (nhiều nhất 1 trên 134 xe). Trong các cấu trúc mạng được thử, chỉ cấu trúc ``cạnh vỉa hè riêng chạy tới vạch dừng, có liên kết đèn riêng và được ưu tiên'' kèm điều kiện kích hoạt là chạy được: không có teleport kẹt, ít nhất 97\% xe đi hết vỉa hè; chi phí cho xe tuân thủ là thời gian đi thẳng tăng từ 42~s lên 59--84~s. Cấu trúc nhập lại kiểu phải nhường đường gây kẹt cứng; thêm làn mép trên cùng cạnh làm dịch mọi chỉ số làn nên không dùng. SUMO không cung cấp chỉ số an toàn cho xe trên vỉa hè.
\item \textbf{3C -- đi ngược chiều (hành vi 8).} Cơ chế vượt bằng làn ngược chiều sẵn có không dùng được: trên đường nhiều làn chỉ xe ưu tiên được dùng; tham số \texttt{lcOpposite} theo tài liệu chỉ áp dụng cho mô hình LC2013, vốn không chạy cùng mô hình làn con; với SL2015 và làn con, xe có phản ứng nhưng đứng lại trên làn ngược chiều (trung vị 54~s). Một làn ảo ngược chiều dành cho \texttt{custom1} chạy được (71/71 xe dùng, 0 va chạm) nhưng SUMO không tính xung đột giữa làn ảo và dòng ngược chiều thật, nên năng lực bị thổi phồng và không có chỉ số an toàn. Lệnh TraCI \texttt{changeSublane} đưa được xe sang làn ngược chiều thật (có tính xung đột) nhưng không kiểm soát được điểm quay về.
\item \textbf{3D -- chờ giữa nút (hành vi 10).} Trên mạng F2, với kịch bản \texttt{target} và 5 seed, số xe dừng lớn nhất trong vùng chờ là 13, 8, 8, 7 và 10, giống kết quả trên vùng chờ 3,5~m ở Bảng~\ref{table:4.box_r21}; không có va chạm nào giữa xe rẽ trái trái phép với các xe khác và không có kẹt cứng.
\end{itemize}

\textbf{Quyết định.} Hành vi (6) và (7) được cài vào mạng đại diện với giá trị giả định lấy từ thiết lập của 3A và 3B (đoạn cọc 30~m; vỉa hè ở nhánh Lê Văn Lương phía tây, rộng 3,0~m, tốc độ 5~m/s, bắt đầu 170~m trước vạch dừng), cùng lựa chọn \texttt{right\_lane} và \texttt{sidewalk} của luật \texttt{queue\_space\_seeking}; trong tệp cấu hình hiện trường đo thực, các trường này để trống. Hành vi (8) được \emph{hoãn}: không cài vào mô hình, chỉ số tương ứng (M11) và lựa chọn \texttt{opposite} của luật không được hiện thực, và hành vi này được ghi vào phần hạn chế. Với hành vi (10), cơ chế điểm chờ nội bộ sẵn có tiếp tục được dùng.

\textbf{Kiểm tra chức năng sau khi cài đặt.} Trên PROXY\_D1, 3 seed, 1800~s, so sánh B1 với mọi xác suất bằng 0 và một biến thể với xác suất chọn làn rẽ phải 0,3, chọn vỉa hè 0,3 và $q = 30$~m: lựa chọn \texttt{right\_lane} được SUMO thực thi ở 93/100, 132/143 và 104/109 lần quyết định; lựa chọn \texttt{sidewalk} chỉ thực thi ở 8/17, 16/33 và 3/25 lần, chủ yếu vì xe ở làn 1--2 không kịp sang lane 0 trong khoảng 27~m trước điểm bắt đầu vỉa hè. Va chạm trong nút không có xu hướng theo luật; số va chạm liên quan can thiệp là 3--11 so với 1--5 của B1.

\textbf{Giả thuyết cần kiểm chứng.} Mức trễ của xe rẽ phải và tỷ lệ dùng vỉa hè phụ thuộc trực tiếp vào chiều dài đoạn cọc, vị trí điểm lên vỉa hè và điểm nhập lại, hiện đều là giả định; cần đo các đại lượng này tại nút trước khi diễn giải các chỉ số M9, M10.

\subsection{Ngưỡng va chạm đóng băng trên mạng F2}
\label{subsection:r28_f2}

\textbf{Thiết lập.} Trước vòng LLM đầu tiên, B0 được chạy trên mạng cuối (F2) với PROXY\_D1--PROXY\_D3, mỗi ngày 3 seed, 1800~s, kiểm tra va chạm trong nút bật, tắt SSM. Hai ngưỡng được tính bằng $\max(2, \lceil p_{90} \rceil)$ của 9 tỷ lệ va chạm mỗi giờ, đóng băng trong cấu hình thí nghiệm và không thay đổi sau khi thấy kết quả của các vòng tối ưu. Kết quả ở Bảng~\ref{table:4.r28}.

\begin{table}[H]
\centering
\small
\begin{tabular}{|p{1.6cm}|p{1.8cm}|p{2.0cm}|p{3.4cm}|p{2.2cm}|}
\hline
\textbf{Ngày / seed} & \textbf{Va chạm làn} & \textbf{Va chạm trong nút} & \textbf{Teleport (tổng / không do va chạm)} & \textbf{Xe chưa chèn cuối} \\ \hline
D1 / 1 & 2 & 41 & 59 / 16 & 339 \\ \hline
D1 / 2 & 0 & 53 & 79 / 26 & 101 \\ \hline
D1 / 3 & 2 & 54 & 56 / 0 & 0 \\ \hline
D2 / 1 & 0 & 47 & 47 / 0 & 0 \\ \hline
D2 / 2 & 0 & 69 & 69 / 0 & 0 \\ \hline
D2 / 3 & 1 & 63 & 91 / 27 & 51 \\ \hline
D3 / 1 & 1 & 78 & 79 / 0 & 0 \\ \hline
D3 / 2 & 0 & 63 & 78 / 15 & 0 \\ \hline
D3 / 3 & 0 & 53 & 60 / 7 & 0 \\ \hline
\end{tabular}
\caption{B0 trên mạng F2, PROXY\_D1--PROXY\_D3 (1800~s mỗi lần chạy, kiểm tra va chạm trong nút bật); số va chạm là số đếm trong một lần chạy}
\label{table:4.r28}
\end{table}

\textbf{Quan sát.} Quy về mỗi giờ, va chạm trên làn là 4, 0, 4, 0, 0, 2, 2, 0, 0 (trung bình 1,333, $p_{90} = 4{,}0$), còn va chạm trong nút nằm trong khoảng 82--156 ($p_{90} = 141{,}6$). Các ngưỡng đóng băng là \textbf{4 va chạm trên làn/h} và \textbf{142 va chạm trong nút/h}; cổng tạm thời ``B0 không quá 5 va chạm trên làn/h'' đạt. Không lần chạy B0 nào đủ điều kiện: các lần chạy bị loại vì teleport không do va chạm, xe chưa chèn được khi kết thúc hoặc va chạm trong nút.

\textbf{Diễn giải.} Với hai làn đi thẳng mỗi nhánh, mạng đại diện bị quá tải ở mức nhu cầu đại diện 8180 xe/h; trong kiểm tra chức năng của nhóm F2, số xe chưa chèn được cuối lần chạy là 300--1200 xe ở cả B0 và B1. Đồ án giữ nguyên nhu cầu và yêu cầu không còn xe chưa chèn khi kết thúc, không sửa mạng để giảm tải; khi không có đặc tả nào đủ điều kiện, đặc tả tốt nhất được chọn theo số teleport kẹt mỗi giờ rồi theo $S$ (Mục~\ref{section:3.8}). Đây là một giới hạn của mô hình nền và phải được nêu cùng mọi kết quả hiệu chỉnh. Ngưỡng 142 va chạm trong nút/h gần như không lọc được lần chạy nào; nó chủ yếu ghi nhận rằng va chạm trong nút của mô hình làn con là rất phổ biến ngay cả ở cấu hình mặc định.

\subsection{Kết quả tạm thời của các chiến dịch hiệu chỉnh}
\label{subsection:interim}

\textbf{Phạm vi và điều kiện.} Mục này báo cáo trạng thái các chiến dịch tại ngày 27/10/2026, khi chúng \emph{chưa hoàn thành}. Chiến dịch B2 (Optuna trên cấu trúc B1) đã dùng hết 36/36 lần chạy, tức 12 thử nghiệm. Hai chiến dịch P1 (LLM chỉnh tham số trên cấu trúc B1) và P2 (LLM sinh cả cấu trúc) mới hoàn thành 2 trong 3 vòng theo kế hoạch, tức 8 đặc tả với 24/36 lần chạy đã có kết quả. Vòng 3 đã được sinh đặc tả và đã bắt đầu chạy trước khi bị tạm dừng theo yêu cầu của người thực hiện; vì vậy bộ đếm ngân sách đã ghi 33 lần chạy (P1) và 27 lần chạy (P2), nhưng chưa có kết quả nào của vòng 3 được lưu và vòng này không được báo cáo ở đây. \textbf{Mọi đặc tả P1 và P2 do mô hình thay thế Claude Sonnet sinh} (\texttt{provenance.model} \texttt{claude-sonnet (stand-in)}), không phải LLM cục bộ, nên các số dưới đây không phải kết quả của phương pháp đề xuất. Mỗi đặc tả được đánh giá bằng 3 lần chạy 1800~s (PROXY\_D1, D2, D3, seed 1), trên dữ liệu đại diện \texttt{LITERATURE\_PROXY}. Để có mốc so sánh, B0 và B1 được chấm điểm theo đúng cách này trong một chiến dịch tham chiếu riêng (REF). Bảng~\ref{table:4.interim_summary} tóm tắt theo chiến dịch và Bảng~\ref{table:4.interim_specs} liệt kê từng đặc tả. Các tỷ lệ mỗi giờ là trung bình của 3 lần chạy; vì mỗi lần chạy dài 0,5~h, một sự kiện tương ứng 2 sự kiện/h trong một lần chạy.

\begin{table}[H]
\centering
\small\setlength{\tabcolsep}{3pt}
\begin{tabular}{|p{1.3cm}|p{1.3cm}|p{1.8cm}|p{2.0cm}|p{1.4cm}|p{1.5cm}|p{1.4cm}|p{2.4cm}|}
\hline
\textbf{Chiến dịch} & \textbf{Đặc tả (lần chạy)} & \textbf{Hợp lệ (lần trả lời / mẫu)} & \textbf{$S$ nhỏ nhất / trung vị / lớn nhất} & \textbf{Va chạm làn/h} & \textbf{Va chạm tổng/h} & \textbf{Teleport kẹt/h} & \textbf{Tốt nhất (quy tắc dự phòng)} \\ \hline
B2 & 12 (36) & -- & 1,182 / 1,329 / 1,523 & 2,0--46,0 & 90,7--168,0 & 12,0--76,7 & \texttt{B2\_\allowbreak t010}: 12,0/h; $S$ 1,374 \\ \hline
P1 (thay thế) & 8 (24) & 8/8 / 8/8 & 1,159 / 1,188 / 1,262 & 0,0--1,3 & 93,3--146,0 & 24,0--54,0 & \texttt{P1\_\allowbreak it02\_\allowbreak s00}: 24,0/h; $S$ 1,167 \\ \hline
P2 (thay thế) & 8 (24) & 8/11 / 8/8 & 1,047 / 1,201 / 1,299 & 6,0--21,3 & 136,7--186,0 & 7,3--30,0 & \texttt{P2\_\allowbreak it02\_\allowbreak s03}: 7,3/h; $S$ 1,084 \\ \hline
B0 (REF) & 1 (3) & -- & 1,181 & 2,0 & 112,7 & 10,7 & -- \\ \hline
B1 (REF) & 1 (3) & -- & 1,216 & 0,0 & 107,3 & 45,3 & -- \\ \hline
\end{tabular}
\caption{Tóm tắt tạm thời theo chiến dịch (PROXY\_D1--D3, seed 1, 1800~s; mọi đặc tả đều bị loại bởi tiêu chí đủ điều kiện; ``Teleport kẹt/h'' gồm teleport do kẹt và do kẹt khi nhường đường; ``Tốt nhất'' theo quy tắc dự phòng: ít teleport kẹt/h nhất, rồi $S$ nhỏ nhất)}
\label{table:4.interim_summary}
\end{table}

\begin{table}[H]
\centering
\footnotesize\setlength{\tabcolsep}{3pt}
\begin{tabular}{|p{1.2cm}|p{2.6cm}|p{1.5cm}|p{1.6cm}|p{1.6cm}|p{1.3cm}|p{1.6cm}|p{1.6cm}|}
\hline
\textbf{Chiến dịch} & \textbf{Đặc tả} & \textbf{Va chạm làn} & \textbf{Va chạm tổng} & \textbf{Teleport kẹt} & \textbf{$S$} & \textbf{$S_{\text{sourced}}$} & \textbf{$S_{\text{assumed}}$} \\ \hline
B2 & \texttt{B2\_t000} & 29,3 & 168,0 & 27,3 & 1,302 & 0,547 & 0,754 \\ \hline
B2 & \texttt{B2\_t001} & 7,3 & 128,0 & 60,0 & 1,356 & 0,531 & 0,825 \\ \hline
B2 & \texttt{B2\_t002} & 9,3 & 113,3 & 43,3 & 1,286 & 0,519 & 0,768 \\ \hline
B2 & \texttt{B2\_t003} & 2,0 & 96,7 & 20,0 & 1,233 & 0,511 & 0,722 \\ \hline
B2 & \texttt{B2\_t004} & 14,0 & 115,3 & 43,3 & 1,523 & 0,522 & 1,002 \\ \hline
B2 & \texttt{B2\_t005} & 7,3 & 112,0 & 13,3 & 1,228 & 0,513 & 0,715 \\ \hline
B2 & \texttt{B2\_t006} & 12,7 & 110,0 & 50,0 & 1,428 & 0,526 & 0,902 \\ \hline
B2 & \texttt{B2\_t007} & 2,7 & 93,3 & 38,0 & 1,182 & 0,490 & 0,692 \\ \hline
B2 & \texttt{B2\_t008} & 21,3 & 120,0 & 49,3 & 1,450 & 0,502 & 0,948 \\ \hline
B2 & \texttt{B2\_t009} & 9,3 & 100,0 & 31,3 & 1,361 & 0,516 & 0,844 \\ \hline
B2 & \texttt{B2\_t010} & 46,0 & 166,0 & 12,0 & 1,374 & 0,526 & 0,848 \\ \hline
B2 & \texttt{B2\_t011} & 3,3 & 90,7 & 76,7 & 1,262 & 0,533 & 0,729 \\ \hline
P1 & \texttt{P1\_it01\_s00} & 0,0 & 93,3 & 25,3 & 1,199 & 0,541 & 0,658 \\ \hline
P1 & \texttt{P1\_it01\_s01} & 1,3 & 102,0 & 37,3 & 1,262 & 0,532 & 0,730 \\ \hline
P1 & \texttt{P1\_it01\_s02} & 0,7 & 115,3 & 28,7 & 1,192 & 0,513 & 0,680 \\ \hline
P1 & \texttt{P1\_it01\_s03} & 0,0 & 100,0 & 36,7 & 1,184 & 0,543 & 0,642 \\ \hline
P1 & \texttt{P1\_it02\_s00} & 0,7 & 140,7 & 24,0 & 1,167 & 0,506 & 0,661 \\ \hline
P1 & \texttt{P1\_it02\_s01} & 0,7 & 146,0 & 54,0 & 1,172 & 0,506 & 0,666 \\ \hline
P1 & \texttt{P1\_it02\_s02} & 0,0 & 118,0 & 45,3 & 1,159 & 0,512 & 0,647 \\ \hline
P1 & \texttt{P1\_it02\_s03} & 0,0 & 120,7 & 49,3 & 1,241 & 0,519 & 0,722 \\ \hline
P2 & \texttt{P2\_it01\_s00} & 19,3 & 186,0 & 21,3 & 1,207 & 0,463 & 0,744 \\ \hline
P2 & \texttt{P2\_it01\_s01} & 12,7 & 173,3 & 18,0 & 1,219 & 0,493 & 0,726 \\ \hline
P2 & \texttt{P2\_it01\_s02} & 8,0 & 176,0 & 21,3 & 1,299 & 0,470 & 0,829 \\ \hline
P2 & \texttt{P2\_it01\_s03} & 9,3 & 162,7 & 24,7 & 1,236 & 0,517 & 0,719 \\ \hline
P2 & \texttt{P2\_it02\_s00} & 6,0 & 136,7 & 14,0 & 1,047 & 0,398 & 0,649 \\ \hline
P2 & \texttt{P2\_it02\_s01} & 21,3 & 168,7 & 30,0 & 1,149 & 0,459 & 0,690 \\ \hline
P2 & \texttt{P2\_it02\_s02} & 11,3 & 153,3 & 24,7 & 1,196 & 0,499 & 0,697 \\ \hline
P2 & \texttt{P2\_it02\_s03} & 9,3 & 152,7 & 7,3 & 1,084 & 0,417 & 0,667 \\ \hline
REF & \texttt{B0} & 2,0 & 112,7 & 10,7 & 1,181 & 0,501 & 0,680 \\ \hline
REF & \texttt{B1\_manual} & 0,0 & 107,3 & 45,3 & 1,216 & 0,525 & 0,691 \\ \hline
\end{tabular}
\caption{Từng đặc tả đã đánh giá (đơn vị: va chạm/h, teleport/h; mọi đặc tả bị loại; P1, P2 do mô hình thay thế sinh). Nguồn: \texttt{report/2026-10-27/data/campaign\_specs.csv} và chiến dịch REF}
\label{table:4.interim_specs}
\end{table}

\textbf{Quan sát.}
\begin{itemize}
\item Không có đặc tả nào của ba chiến dịch đủ điều kiện. Mọi đặc tả đều có teleport kẹt (ít nhất 7,3/h), trong khi tiêu chí đủ điều kiện yêu cầu 0; vì vậy đặc tả tốt nhất của mỗi chiến dịch được chọn bằng quy tắc dự phòng. B0 và B1 chấm theo cùng cách cũng bị loại.
\item Ở B2, đặc tả tốt nhất theo quy tắc dự phòng là \texttt{B2\_t010} (12,0 teleport kẹt/h), dù nó có tỷ lệ va chạm trên làn cao nhất chiến dịch (46,0/h) và $S$ (1,374) cao hơn đặc tả có $S$ nhỏ nhất (\texttt{B2\_t007}, 1,182); quy tắc dự phòng không áp mức trần va chạm. Giá trị $S$ tốt nhất của B2 giảm từ 1,302 ở thử nghiệm đầu xuống 1,182 ở thử nghiệm thứ 8 và không giảm thêm trong 4 thử nghiệm cuối.
\item Tám đặc tả P1 có va chạm trên làn 0,0--1,3/h, thấp hơn khoảng 2,0--46,0/h của 12 đặc tả B2 trên cùng cấu trúc, nhưng có teleport kẹt 24,0--54,0/h, cao hơn đặc tả tốt nhất của B2.
\item P2 có $S$ nhỏ nhất (1,047) và teleport kẹt thấp nhất (7,3/h) trong các đặc tả đã đánh giá, nhưng mọi đặc tả P2 có va chạm trên làn 6,0--21,3/h, vượt ngưỡng 4/h, và va chạm tổng 136,7--186,0/h, cao hơn hai chiến dịch còn lại.
\item B0 chấm theo cùng cách có $S = 1{,}181$ và 10,7 teleport kẹt/h. Chỉ 6/28 đặc tả có $S$ nhỏ hơn B0 (3 của P1, 3 của P2, không có đặc tả nào của B2) và chỉ 1/28 đặc tả (\texttt{P2\_it02\_s03}) có ít teleport kẹt hơn B0. Như vậy, theo chính quy tắc chọn dự phòng (teleport kẹt rồi $S$), chưa đặc tả nào của B2 và P1 tốt hơn B0; chỉ \texttt{P2\_it02\_s03} (do mô hình thay thế sinh; 7,3/h, $S = 1{,}084$) tốt hơn B0 ở cả hai tiêu chí. B1, cấu trúc xuất phát của B2 và P1, có $S = 1{,}216$, 0,0 va chạm trên làn/h và 45,3 teleport kẹt/h; số đặc tả có $S$ nhỏ hơn B1 là 1/12 ở B2, 6/8 ở P1 và 5/8 ở P2, và số đặc tả có ít teleport kẹt hơn B1 là 8/12, 5/8 và 8/8.
\item Về tính hợp lệ: ở P1, cả 8 lần trả lời đều hợp lệ ngay lần đầu. Ở P2, vòng 1 chỉ có 1/4 đặc tả hợp lệ ở lần trả lời đầu; 3 đặc tả còn lại cùng mắc một lỗi logic (khai báo đồng thời \texttt{lcPushy} và \texttt{lcPushyGap}) và đều được sửa ở lần thử thứ hai sau khi nhận thông báo lỗi; vòng 2 có 4/4 hợp lệ ngay lần đầu. Tính theo lần trả lời, P2 có 8/11 hợp lệ; tính theo trạng thái cuối của mẫu, cả hai chiến dịch đạt 8/8.
\end{itemize}

\textbf{Diễn giải có cơ sở.} Trên mạng đại diện quá tải, quy tắc chọn dự phòng ưu tiên giảm ùn tắc chứ không ưu tiên độ tương đồng, nên ``đặc tả tốt nhất'' ở đây có thể có $S$ cao hơn hoặc va chạm nhiều hơn các đặc tả khác cùng chiến dịch, như trường hợp \texttt{B2\_t010}. Các mức $S$ của mọi phương án nằm trong khoảng hẹp (1,05--1,52) và đa số không tốt hơn B0, nên ở giai đoạn này chưa phương án nào cho thấy cải thiện rõ ràng so với cấu hình mặc định. Cơ chế phản hồi lỗi của bộ kiểm tra hợp lệ đã hoạt động đúng như thiết kế trong một trường hợp thực (lỗi \texttt{lcPushy}/\texttt{lcPushyGap} được sửa sau một lần thử).

\textbf{Vì sao chưa thể so sánh.} So sánh giữa các chiến dịch ở đây \emph{không} phải kết quả cuối và không được dùng để kết luận phương án nào tốt hơn, vì: (i) P1 và P2 mới dùng 2/3 ngân sách, trong khi B2 đã dùng hết; (ii) mỗi ngày chỉ có một seed, nên chưa tách được biến động do seed; (iii) mọi đặc tả đều bị loại bởi tiêu chí đủ điều kiện và đặc tả tốt nhất được chọn theo quy tắc dự phòng; (iv) mô hình nền quá tải ở nhu cầu đại diện; (v) P1, P2 dùng mô hình thay thế chứ không phải LLM cục bộ; (vi) dữ liệu mục tiêu là dữ liệu đại diện, và các khóa có mục tiêu giả định chiếm hơn một nửa giá trị $S$ ở mọi đặc tả. Đánh giá E1, E2 trên ngày đánh giá với nhiều seed và kiểm định thống kê (Mục~\ref{subsection:planned_experiments}) là bước cần thiết trước khi đưa ra bất kỳ kết luận so sánh nào.

\subsection{Các thực nghiệm dự kiến}
\label{subsection:planned_experiments}

Các thực nghiệm sau chưa được thực hiện. Trước tiên, hai chiến dịch P1 và P2 cần dùng hết 12 lần chạy còn lại (vòng 3, và thêm vòng nếu còn ngân sách, tối đa 5 vòng), lý tưởng là với LLM cục bộ thay cho mô hình thay thế. Sau đó, phương án H chạy Optuna với cùng ngân sách 36 lần chạy trên cấu trúc của đặc tả tốt nhất của P2. Đánh giá cuối dự kiến chạy B0, B1, B2, P1, P2 và H trên hai ngày đánh giá PROXY\_D4--PROXY\_D5 với ít nhất 20 seed; thí nghiệm E1 so sánh P1 với B2, thí nghiệm E2 so sánh B2, H và P2, so sánh từng cặp bằng kiểm định Mann--Whitney có hiệu chỉnh Holm và kích thước hiệu ứng Cliff's delta, kèm một lần chạy chẩn đoán 3600~s $\times$ 3 seed cho B0 và đặc tả tốt nhất trên PROXY\_D4. Nếu còn thời gian, đồ án dự kiến thực hiện thêm phân tích loại bỏ (ablation) bằng cách tắt lần lượt từng luật TraCI hoặc từng nhóm tham số, và kịch bản chính sách với lưu lượng rẽ trái trái phép bằng 0\%, 50\% và 100\% mức quan sát, đo độ trễ và số xung đột; kịch bản chính sách sẽ được ghi rõ là ngoại suy. Mô hình LLM cục bộ dùng cho kết quả chính thức vẫn chưa được chốt \textbf{[TODO]}.

\section{Triển khai}

Hệ thống được phát triển trên một máy Windows, với SUMO 1.27.0 cài đặt sẵn và biến môi trường \texttt{SUMO\_HOME} trỏ tới thư mục cài đặt, rồi chuyển sang một máy Linux, nơi SUMO 1.27.0 được cài dưới dạng gói Python \texttt{eclipse-sumo} trong môi trường ảo và tệp đường dẫn của môi trường ảo tự đặt \texttt{SUMO\_HOME} \textbf{[TODO: cấu hình phần cứng của hai máy (CPU, RAM, GPU)]}. Tám lần chạy B0 đã có trên Windows được chạy lại trên Linux và cho đúng từng số va chạm và teleport. Các chiến dịch B2, P1, P2 và REF chạy trên máy Linux. Môi trường Python là một môi trường ảo tạo bằng Python 3.12, dựng bằng cùng một script trên cả hai máy. Gói dự án được cài ở chế độ có thể chỉnh sửa kèm nhóm phụ thuộc phát triển; \texttt{libsumo} và \texttt{sumo-data} phiên bản 1.27.0 được cài không kèm phụ thuộc; một tệp đường dẫn trong môi trường ảo đưa thư mục công cụ của SUMO vào đường dẫn tìm mô-đun để nạp \texttt{traci} và \texttt{sumolib}. Toàn bộ kiểm thử chạy bằng \texttt{pytest} từ thư mục gốc của dự án.

Quy trình tái lập gồm các lệnh mô-đun Python theo thứ tự của Bảng~\ref{table:4.2}: \texttt{thesis\_sumo.net.build} với tệp cấu hình hiện trường đại diện để dựng mạng; \texttt{thesis\_sumo.proxy.generate} để sinh năm ngày dữ liệu đại diện và tệp tuyến tương ứng; \texttt{thesis\_sumo.rules.compiler} để biên dịch một đặc tả (bắt buộc truyền tệp cấu hình hiện trường đại diện, vì cấu hình đo thực chưa có bề rộng làn); và \texttt{thesis\_sumo.sim.runner} để chạy mô phỏng, với các tham số mã đặc tả, ngày, danh sách seed, tệp cấu hình hiện trường, thời điểm kết thúc, số tiến trình song song, chế độ SSM và tùy chọn mở \texttt{sumo-gui}. Thí nghiệm vùng chờ giữa nút và các thí nghiệm Phase 3A--3C có script riêng. Các chiến dịch được chạy bằng \texttt{thesis\_sumo.baselines.optuna\_search} (B2, H) và \texttt{thesis\_sumo.llm.loop} (P1 với chế độ \texttt{params\_only} và đặc tả cơ sở B1, P2 với chế độ \texttt{full}), cùng ngân sách 36, cùng tệp cấu hình hiện trường đại diện và cùng mức trần va chạm B0 1,333/h; với mô hình thay thế, lệnh \texttt{loop pending} liệt kê các yêu cầu chưa có phản hồi, và chạy lại cùng lệnh sau khi đã có phản hồi sẽ tiếp tục chiến dịch. Dữ liệu của mỗi chiến dịch nằm trong \texttt{experiments/\allowbreak campaigns/\allowbreak <chiến dịch>/} (ngân sách, đặc tả, đặc tả đã biên dịch, lần chạy, trạng thái vòng lặp hoặc nghiên cứu Optuna) và nhật ký LLM nằm trong \texttt{experiments/\allowbreak llm\_logs/\allowbreak <chiến dịch>/}.

Mỗi lần chạy ghi vào một thư mục riêng theo dạng \texttt{experiments/\allowbreak runs/\allowbreak <spec\_id>\_\_<day>/\allowbreak seed\_XX/}, với \texttt{spec\_id} là mã đặc tả, \texttt{day} là ngày dữ liệu và \texttt{XX} là seed; thư mục của seed bị xóa và tạo lại nếu chạy lại. Thư mục gồm tệp cấu hình \texttt{run.sumocfg} và tệp khai báo đầu ra, các đầu ra của SUMO (thông tin chuyến đi, quỹ đạo nén, chuyển làn, thống kê, va chạm, trạng thái đèn, detector), nhật ký \texttt{sumo.log}, tệp \texttt{vtypes.json} ghi các vType SUMO đã nạp, \texttt{interventions.csv} nếu có luật TraCI, \texttt{manifest.json} với các trường nêu ở Mục~\ref{subsection:package_design}, và thư mục các bảng CSV theo lược đồ quan sát. Các đầu ra thô không được đưa vào kho mã; kho mã chỉ lưu mã, cấu hình, tệp vá, mạng đã sinh, dữ liệu đại diện và các bảng kết quả nhỏ. Thời gian thực của một lần chạy B0 1800~s trên dữ liệu đại diện, khi chạy 3 lần song song, là 88--463~s.

\end{document}
